\documentclass{article}
\usepackage{amsthm}

% \theoremstyle declarations are stripped by build_minimal_preamble.
% Environments defined after \theoremstyle{definition} should have
% upright (non-italic) body text. Without the \theoremstyle, they
% inherit the previous 'plain' style and render body text in italics.
\newtheorem{theorem}{Theorem}
\theoremstyle{definition}
\newtheorem{definition}{Definition}
\newtheorem{remark}{Remark}

\begin{document}
\begin{theorem}
  This is a theorem — body should be italic (plain style).
\end{theorem}
\begin{definition}
  This is a definition — body should NOT be italic (definition style).
\end{definition}
\begin{remark}
  This is a remark — body should NOT be italic (definition style).
\end{remark}
\end{document}
